\documentclass[notes,10pt, aspectratio=169]{beamer}

\usepackage{pgfpages}
% These slides also contain speaker notes. You can print just the slides,
% just the notes, or both, depending on the setting below. Comment out the want
% you want.
% Note visibility is selected after package loading at the end of this preamble.
%\setbeameroption{show only notes} % Only notes
%\setbeameroption{show notes on second screen=right} % Both

\usepackage{helvet}
\usepackage[default]{lato}
\usepackage{array}
\usepackage{eurosym}
\usepackage{setspace}
\usepackage{animate}
\usepackage{multicol}
\usepackage{breqn}
\usepackage{subcaption}
\usepackage{caption}
\captionsetup{font=scriptsize, justification=raggedright,singlelinecheck=false}
\usepackage[export]{adjustbox}
\usepackage[scale=2]{ccicons}
\usepackage{tikz}
\newenvironment{wideitemize}{\itemize\addtolength{\itemsep}{10pt}}{\enditemize}
\usepackage{natbib}

\usepackage{tikz}
\usepackage{verbatim}
\setbeamertemplate{note page}{\pagecolor{yellow!5}\insertnote}
\usepackage{amsmath}
\usepackage{bm}
%\usepackage{mathpazo}
\usepackage{hyperref}
\usepackage{lipsum}
\usepackage{multimedia}
\usepackage{multirow}
\usepackage{graphicx}
\usepackage{dcolumn}
\usepackage{bbm}
\newcolumntype{d}[0]{D{.}{.}{5}}

\usepackage{changepage}
\usepackage{appendixnumberbeamer}
% \newcommand{\beginbackup}{
%   \newcounter{framenumbervorappendix}
%   \setcounter{framenumbervorappendix}{\value{framenumber}}
%   \setbeamertemplate{footline}
%   {
%      \leavevmode%
%      \hline
%      box{%
%       \begin{beamercolorbox}[wd=\paperwidth,ht=2.25ex,dp=1ex,right]{footlinecolor}%
% %         \insertframenumber  \hspace*{2ex}
%       \end{beamercolorbox}}%
%      \vskip0pt%
%   }
%  }

 \setbeamertemplate{footline}{
    \hbox{%
    \begin{beamercolorbox}[wd=\paperwidth,ht=3ex,dp=1.5ex,leftskip=2ex,rightskip=2ex]{page footer}%
        %\usebeamerfont{title in head/foot}%
        %\insertshorttitle \hfill
        %\insertsection \hfill
        %\tableofcontents[currentsection] \hfill
        \insertsectionnavigationhorizontal{0.7\paperwidth}{}{} \hfill
        \insertframenumber{} / \inserttotalframenumber
    \end{beamercolorbox}}%
}
\newcommand{\backupend}{
   \addtocounter{framenumbervorappendix}{-\value{framenumber}}
   \addtocounter{framenumber}{\value{framenumbervorappendix}}
}


\usepackage{graphicx}
\usepackage[space]{grffile}
\usepackage{booktabs}

% These are my colors -- there are many like them, but these ones are mine.
\definecolor{blue}{RGB}{0,114,178}
\definecolor{red}{RGB}{213,94,0}
\definecolor{yellow}{RGB}{240,228,66}
\definecolor{green}{RGB}{0,158,115}

\setbeamercolor{section in head/foot}{fg=blue}

\useinnertheme[shadow]{rounded}
\setbeamercolor{block title example}{fg=black,bg=orange!50!white}
\setbeamercolor{block body example}{fg=black, bg=orange!30!white}


\hypersetup{
  colorlinks=false,
  linkbordercolor = {white},
  linkcolor = {blue}
}


%% I use a beige off white for my background
\definecolor{MyBackground}{RGB}{255,253,218}

%% Uncomment this if you want to change the background color to something else
%\setbeamercolor{background canvas}{bg=MyBackground}

%% Change the bg color to adjust your transition slide background color!
\newenvironment{transitionframe}{
  \setbeamercolor{background canvas}{bg=yellow}
  \begin{frame}}{
    \end{frame}
}

\setbeamercolor{frametitle}{fg=blue}
\setbeamercolor{title}{fg=black}
%\setbeamertemplate{footline}[frame number]
\setbeamertemplate{navigation symbols}{}
\setbeamertemplate{itemize items}{-}
\setbeamercolor{itemize item}{fg=blue}
\setbeamercolor{itemize subitem}{fg=blue}
\setbeamercolor{enumerate item}{fg=blue}
\setbeamercolor{enumerate subitem}{fg=blue}
\setbeamercolor{button}{bg=white,fg=blue,}



% If you like road maps, rather than having clutter at the top, have a roadmap show up at the end of each section
% (and after your introduction)
% Uncomment this is if you want the roadmap!
% \AtBeginSection[]
% {
%    \begin{frame}
%        \frametitle{Roadmap of Talk}
%        \tableofcontents[currentsection]
%    \end{frame}
% }
\setbeamercolor{section in toc}{fg=blue}
\setbeamercolor{subsection in toc}{fg=red}
%\setbeamersize{text margin left=1em,text margin right=1em}
\setbeamersize{text margin left=0.7cm,text margin right=0.7cm}


\title[]{\textcolor{blue}{Seminar 9 v2:\\Instruments for the Problems in Seminar 8}}
\subtitle{Four problems and worked solutions}
\author{25117 - Econometrics}
\institute[shortinst]{Universitat Pompeu Fabra}

%\date{September 25th, 2024}


% Each problem is followed by worked solution slides. Set \showsolutionsfalse
% below to compile a problem-only deck. All data are hypothetical;
% no external assets are needed. All numerical moments in this version
% are population moments. IV ratios illustrate identification, not exact
% finite-sample estimates. Every regression includes an intercept.
\newif\ifshowsolutions
\showsolutionstrue
\newcommand{\E}{\mathbb{E}}
\newcommand{\Var}{\operatorname{Var}}
\newcommand{\Cov}{\operatorname{Cov}}
\newcommand{\Corr}{\operatorname{Corr}}
\newcommand{\answerpart}[1]{\textcolor{blue}{\textbf{#1.}}\,}

\setbeameroption{hide notes} % Speaker notes hidden; worked solutions are visible.


\date{}
\begin{document}
\begin{frame}[plain]
\maketitle
\end{frame}

\section{Omitted variables}
\begin{frame}{Problem 1: An instrument for class size}
\small
\setlength{\abovedisplayskip}{4pt}
\setlength{\belowdisplayskip}{4pt}
\textbf{Seminar 8, Problem 1 (slides 2--4):} the estimated class-size slope was $-2$ without prior performance $W$ and $-1$ with it. Can IV recover the slope when $W$ is unobserved? Consider a new population:
\[
Y=80-X+W+R.
\]
Here $X$ is class size, $W$ is prior reading performance, and $Y$ is later performance; all are continuous school averages. $W$ is unobserved. Assume
\[
\E[R]=0,\text{ }\Cov(X,R)=0,\text{ and }\frac{\Cov(X,W)}{\Var(X)}=-1.
\]
A continuously varying, randomly assigned capacity allocation $Z$ changes class size. Assume it affects scores only through $X$ and is uncorrelated with $W$ and $R$. Population moments are
\[
\Var(Z)=4,\text{ }\Cov(Z,X)=2,\text{ and }\Cov(Z,Y)=-2.
\]
\begin{enumerate}
\item Calculate the population OLS slope of $Y$ on $X$ with an intercept.
\item State the relevance and exogeneity conditions for $Z$. Why could $Z$ satisfy them even though $X$ does not?
\item Calculate the first-stage slope, reduced-form slope, and IV ratio. Interpret the IV slope and compare it with OLS.
\end{enumerate}
\end{frame}

\ifshowsolutions
\begin{frame}{Solution 1: Isolate variation unrelated to prior performance}
\small
\setlength{\abovedisplayskip}{4pt}
\setlength{\belowdisplayskip}{4pt}
\answerpart{1}Write $Y=80-X+U$, where $U=W+R$. Then
\[
\Cov(X,U)=\Cov(X,W)+\Cov(X,R)=\Cov(X,W).
\]
The population OLS slope is
\[
b_{\mathrm{OLS}}=-1+\frac{\Cov(X,U)}{\Var(X)}=-1-1=\boxed{-2}.
\]
\answerpart{2}Relevance holds because $\Cov(Z,X)=2\ne0$. For exogeneity, use $U=W+R$ and covariance linearity:
\[
\begin{aligned}
\Cov(Z,U)&=\Cov(Z,W+R)\\
&=\Cov(Z,W)+\Cov(Z,R)=0+0=0.
\end{aligned}
\]
The two zeros follow from the assumptions that $Z$ is uncorrelated with both $W$ and $R$.
\par\smallskip
\answerpart{3}The first-stage and reduced-form slopes are
\[
\pi_1=\frac{\Cov(Z,X)}{\Var(Z)}=\frac24=0.5,
\text{ and }\gamma_1=\frac{\Cov(Z,Y)}{\Var(Z)}=\frac{-2}{4}=-0.5.
\]
\[
b_{\mathrm{IV}}=\frac{\gamma_1}{\pi_1}=\frac{-0.5}{0.5}=\boxed{-1}.
\]
Under these assumptions, one more pupil per class lowers the later score by 1 point. IV removes the omitted-variable component that makes the OLS slope $-2$.
\end{frame}
\fi

\section{Measurement error}
\begin{frame}{Problem 2: An instrument for inaccurately recorded hours}
\small
\setlength{\abovedisplayskip}{4pt}
\setlength{\belowdisplayskip}{4pt}
\textbf{Seminar 8, Problem 2 (slides 5--7):} measurement error made the OLS slope converge to $2$ although the true slope was $3$. Can IV recover $3$? Reuse the model with continuous practice hours and scores:
\[
Y=40+3X+U,\text{ and }\widetilde X=X+V.
\]
Assume $\Cov(X,U)=\Cov(X,V)=\Cov(U,V)=0$, $\Var(X)=4$, and $\Var(V)=2$. Errors have mean zero.
\par\smallskip
A continuously varying, randomly assigned booking allowance $Z$ changes true practice hours. Assume it affects scores only through true hours and is unrelated to recording errors:
\[
\Var(Z)=1,\text{ }\Cov(Z,X)=1.5,\text{ and }\Cov(Z,U)=\Cov(Z,V)=0.
\]
\begin{enumerate}
\item Rewrite the outcome equation using $\widetilde X$ and identify its composite error. Calculate the population OLS slope using recorded hours.
\item Show that $Z$ is relevant for $\widetilde X$ and uncorrelated with the composite error. Calculate $\Cov(Z,Y)$.
\item Compute the IV ratio using the observed regressor $\widetilde X$. Describe both regression stages. Would validity still follow if $Z$ were correlated with recording error $V$?
\end{enumerate}
\end{frame}

\ifshowsolutions
\begin{frame}{Solution 2: Instrument the recorded regressor}
\small
\setlength{\abovedisplayskip}{4pt}
\setlength{\belowdisplayskip}{4pt}
\answerpart{1}Substitute $X=\widetilde X-V$:
\[
Y=40+3\widetilde X+\widetilde U,\text{ where }\widetilde U=U-3V.
\]
\[
\Cov(\widetilde X,\widetilde U)=-3\Var(V)=-6,
\text{ and }b_{\mathrm{OLS}}=3\frac4{4+2}=\boxed2.
\]
\answerpart{2}The instrument isolates true-hour variation rather than recording noise:
\[
\begin{aligned}
\Cov(Z,\widetilde X)&=\Cov(Z,X)+\Cov(Z,V)=1.5,\\
\Cov(Z,\widetilde U)&=\Cov(Z,U)-3\Cov(Z,V)=0,\\
\Cov(Z,Y)&=3\Cov(Z,X)+\Cov(Z,U)=4.5.
\end{aligned}
\]
\answerpart{3}Consequently,
\[
b_{\mathrm{IV}}=\frac{\Cov(Z,Y)}{\Cov(Z,\widetilde X)}
=\frac{4.5}{1.5}=\boxed{3\text{ points per true practice hour}}.
\]
First regress recorded hours $\widetilde X$ on an intercept and $Z$. Then regress $Y$ on an intercept and the fitted recorded hours. The population first-stage slope is 1.5, the reduced-form slope is 4.5, and their ratio is 3.
\par\smallskip
If $\Cov(Z,V)\ne0$, then $\Cov(Z,\widetilde U)=-3\Cov(Z,V)\ne0$: the instrument would be invalid under the other maintained assumptions.
\end{frame}
\fi

\section{Selection}
\begin{frame}{Problem 3: Does an instrument survive sample selection?}
\small
\setlength{\abovedisplayskip}{4pt}
\setlength{\belowdisplayskip}{4pt}
\textbf{Seminar 8, Problem 3 (slides 8--9):} retaining $Y\geq3$ gave conditional error means $1,1,0$ at $X=0,1,2$, breaking zero conditional mean. Does IV fix this? Consider a continuous-variable version:
\[
Y=2+X+U.
\]
A continuous instrument $Z$ is valid in the full population: $\Cov(Z,U)=0$ and $\Cov(Z,X)\ne0$. Only observations with $Y\geq3$ are retained; write $K=1$ for retention. Suppose the selected-population moments are
\[
\Cov(Z,X\mid K=1)=0.60,\text{ and }\Cov(Z,U\mid K=1)=-0.20.
\]
\begin{enumerate}
\item Derive the selected-population IV ratio. Calculate it and compare it with the structural slope 1.
\item Explain why full-population validity does not establish validity after selection on $Y$.
\item A redesigned study selects participants using \emph{pre-training} performance, then randomly assigns a continuous encouragement $Q$. Assume $Q$ affects current $Y$ only through $X$ and has no spillovers. Among these participants, $\Cov(Q,X)=0.50$ and $\Cov(Q,Y)=0.50$. Calculate IV and explain what changed in the design.
\end{enumerate}
\end{frame}

\ifshowsolutions
\begin{frame}{Solution 3: Validity must hold in the sample used}
\small
\setlength{\abovedisplayskip}{4pt}
\setlength{\belowdisplayskip}{4pt}
\answerpart{1}Take covariances within the selected population:
\[
\Cov(Z,Y\mid K=1)=\Cov(Z,X\mid K=1)+\Cov(Z,U\mid K=1)=0.40.
\]
\[
b_{\mathrm{IV},K=1}
=1+\frac{\Cov(Z,U\mid K=1)}{\Cov(Z,X\mid K=1)}
=1-\frac{0.20}{0.60}=\boxed{\frac23\ne1}.
\]
\answerpart{2}Selection depends on both $X$ and $U$, because $Y\geq3$ means $U\geq1-X$. Since $Z$ shifts $X$, keeping only high outcomes can make $Z$ correlated with $U$ among retained observations. Here relevance survives, but exogeneity fails.
\par\smallskip
\answerpart{3}In the redesigned study, selection occurs before random assignment. Under the stated exclusion assumption, assignment supports instrument exogeneity within the enrolled population. Thus
\[
b_{\mathrm{IV},\mathrm{new}}=\frac{\Cov(Q,Y)}{\Cov(Q,X)}
=\frac{0.50}{0.50}=\boxed1.
\]
This identifies the constant slope for enrolled participants. It relies on a new design; applying IV to the original outcome-selected records does not repair their selection bias automatically.
\end{frame}
\fi

\section{Simultaneity}
\begin{frame}{Problem 4: An instrument that shifts the regressor}
\small
\setlength{\abovedisplayskip}{4pt}
\setlength{\belowdisplayskip}{4pt}
\textbf{Seminar 8, Problem 4 (slides 10--11):} feedback gave $\Cov(X,U)=2$ and an OLS slope $2+2/\Var(X)$ instead of $2$. Can IV isolate the structural slope?
\par\smallskip
Assume we can find a new observed variable $Z$ that affects $X$ but has no direct effect on $Y$. Keep the original unobserved shock $V$ and add $Z$:
\[
Y=2X+U,\text{ and }X=0.25Y+Z+V.
\]
All variables are continuous. The variables $Z$, $U$, and $V$ are mutually independent, have mean zero, and satisfy
\[
\Var(Z)=1,\text{ }\Var(U)=4,\text{ and }\Var(V)=1.
\]
$Z$ can affect $Y$ indirectly through $X$.
\begin{enumerate}
\item Solve for $X$ in terms of $Z$, $U$, and $V$. Calculate $\Cov(X,U)$, $\Var(X)$, and the population OLS slope of $Y$ on $X$.
\item Calculate $\Cov(Z,X)$ and $\Cov(Z,Y)$. Explain why $Z$ can be a valid instrument although $X$ is endogenous.
\item Calculate the IV slope. Write the population first-stage and reduced-form fitted lines, and show how two stages recover the structural slope.
\end{enumerate}
\end{frame}

\ifshowsolutions
\begin{frame}{Solution 4: Use the external shift rather than the feedback}
\small
\setlength{\abovedisplayskip}{3pt}
\setlength{\belowdisplayskip}{3pt}
\answerpart{1}Substitution gives $X=0.5X+0.25U+Z+V$, hence
\[
X=0.5U+2Z+2V.
\]
\[
\Cov(X,U)=0.5(4)=2,\text{ and }\Var(X)=0.25(4)+4(1)+4(1)=9.
\]
\[
b_{\mathrm{OLS}}=2+\frac{\Cov(X,U)}{\Var(X)}=2+\frac29=\boxed{\frac{20}{9}}.
\]
\answerpart{2}Independence and exclusion imply
\[
\Cov(Z,U)=0,\text{ }\Cov(Z,X)=2\Var(Z)=2,
\]
\[
\Cov(Z,Y)=2\Cov(Z,X)+\Cov(Z,U)=4.
\]
The instrument shifts $X$ without carrying the outcome shock $U$.
\par\smallskip
\answerpart{3}The IV ratio is $b_{\mathrm{IV}}=4/2=\boxed2$. Since all variables have mean zero, the population fitted lines have zero intercepts:
\[
X^{\mathrm{proj}}=2Z,\text{ and }Y^{\mathrm{proj}}=4Z.
\]
Regressing $Y$ on the first-stage projection gives
\[
\frac{\Cov(2Z,Y)}{\Var(2Z)}=\frac{2(4)}{4(1)}=\boxed2.
\]
\end{frame}
\fi

\end{document}
